The two benchmarks give different roles to a kernel fitted after neural training. On structural mechanics, the neural members and their combination account for most of the accuracy, with a smaller repeatable gain from residual correction. On OCO-2, the representation learned by the network is the main advantage: the direct-target kernel head performs substantially better on frozen features than on raw inputs. Coordinate selection then balances errors across two metrics that weight the reduced outputs very differently.

The structural-mechanics score of $\errSix\pm\errSixSd\%$ is close to PARA-Net's published $\pubBestHi\%$; the two studies differ in training subset and seed protocol, so the comparison is indicative rather than a ranking. The improvement in the lower-data regime is larger, subject to uncertainty about the source study's exact training subset. The paired reruns in Section~\ref{sec:implementation-sensitivity} change target centering while repeating downstream training and fitting. Their mean change is only $\sensCenterDelta$ percentage points, with changes in both directions across seeds. Replacing correction labels by residuals of the query-time mean also has a small effect under the fixed historical kernel. These checks quantify two implementation choices on this benchmark.

The ensemble analysis gives a more specific conclusion. For the sixty trained predictors, the empirical lower bound of $\saBound\%$ lies close to the hindsight-optimal global convex RMS error of $\saSixtyOracleRms\%$. There is little room left for this form of averaging within the pool. Both individual accuracy and residual alignment matter: completing the FNO schedule improves its own error without improving the five-member pipeline, whereas adding the UNet improves the corrected pipeline by $\unetGain$ percentage points. That last result comes from an ablation of the full estimator, whose per-pixel weights and kernel correction extend beyond the class covered by the convex bound. The shared residual is therefore evidence about these predictors, not an irreducible floor for the benchmark.

The OCO-2 comparison benefits from access to the source emulator's predictions on the same test states. Its scope is the retained forty-component representation: full-spectrum truncation, instrument response, Jacobians, and retrieval accuracy are not evaluated. The coordinate-selected predictor stays close to the best head on the reduced criterion while improving radiance error, rather than dominating every constituent on both. A frozen linear-readout control supports the benefit of the nonlinear kernel head in these experiments. The wider matched kernel search in Section~\ref{sec:oco-grid-sensitivity} preserves the feature advantage over raw-input kernels, but its validation-selected coordinate combination has higher mean reduced error in all three bands. A larger candidate set does not guarantee lower test error.

The theory helps interpret these results without turning the diagnostics into error bars. The native-space bound separates a residual norm from a design factor and gives the exact worst-case cost of a positive ridge parameter. Applying it to the implemented correction also requires the label-mismatch term in Section~\ref{sec:theory-or}. The bound is conditional on the residual's RKHS membership and norm, which are not assumed for the benchmark target. Fitted norms describe the finite-design regressions, and their ratios do not determine ratios of the unknown target norms. The pullback identity likewise characterizes the space induced by neural features without guaranteeing that its norm is smaller than the raw-input norm.

For uncertainty quantification, the simpler scale works better here. Both kernel power functions and ensemble disagreement rank errors weakly, and a constant-radius calibrated set is narrower on average. All three constructions have the same marginal split-conformal guarantee when the predictor and scale are fixed independently of exchangeable calibration and evaluation data. The exact Beta and Beta--binomial reference laws further require i.i.d. continuous scores.

The supplement provides the proofs, the paired sensitivity results, and the computational details behind the reported numbers.
